\bookchapter{Глава 59. Знакомство с iOS — экосистема Apple и экосистема разработки}

\chapterintro{Поздравляем с этим решением — перейти от «человека, который пользуется iPhone» к «человеку, который пишет приложения для iPhone». Прежде чем браться за дело, разложим карту: откуда взялся iOS? Что значит для разработчика «сад с высокими стенами» Apple? Что происходит между моментом, когда ваши пальцы стучат по клавиатуре, и появлением приложения на телефоне пользователя?}

\booksection{Происхождение iOS: один телефон, изменивший целую отрасль}

В январе 2007 года Стив Джобс на презентации достал из кармана первый iPhone. Тогда он даже не поддерживал сторонние приложения — и только в 2008 году Apple выпустила App Store и iPhone SDK, разрешив разработчикам со всего мира писать программы для iPhone. В тот момент родилась «экономика приложений»: сегодня в App Store миллионы приложений, а суммарные выплаты разработчикам превысили сотни миллиардов долларов.

Сама система тоже эволюционировала: сначала она называлась iPhone OS, в 2010 году сменила имя на iOS; номера версий дошли до iOS 17 (2023), а в 2025 году Apple перевела именование на «годовой» принцип — iOS 26 вышла в 2025 году, а текущая версия в августе 2026-го — iOS 27. Названия меняются, но одно не менялось никогда: привычка пользователей iOS платить и репутация качества приложений по праву считаются лучшими в мобильной экосистеме — одно и то же приложение на iOS часто приносит заметно больший доход, чем на Android.

\begin{center}
\begin{tabularx}{\linewidth}{|X|X|X|}
\hline
\textbf{Событие} & \textbf{Год} & \textbf{Что важно для разработчиков} \tabularnewline
\hline
Запуск App Store & 2008 & Открытие экосистемы сторонней разработки \tabularnewline
Выпуск Swift & 2014 & Появление современного языка разработки для iOS \tabularnewline
Выпуск SwiftUI & 2019 & Декларативный фреймворк интерфейсов, основная линия этой книги \tabularnewline
iOS 26 (первый год нового именования) & 2025 & Названия привязаны к году; повсеместная интеграция Apple Intelligence \tabularnewline
iOS 27 & 2026 & Текущая новейшая версия (ориентир этой книги) \tabularnewline
\hline
\end{tabularx}
\end{center}

\begin{center}\small\color{gray}Ключевые вехи разработки iOS (по август 2026 года)\end{center}

\booksection{Сад с высокими стенами: «правила» и «дивиденды» экосистемы Apple}

Экосистему Apple часто описывают как «сад с высокими стенами» (Walled Garden) — правил много, контроль строгий: приложение должно пройти проверку, чтобы попасть в магазин; разработка возможна только на Mac и в Xcode; распространение идёт в основном через App Store. Звучит тесновато? Но для разработчика эти правила приносят вполне осязаемые дивиденды:

\begin{itemize}
  \item Крайне малая фрагментация устройств: моделей iPhone всего несколько, размеры экранов и версии системы очень однородны — адаптировав 5 моделей, вы охватите 95\% пользователей (на Android может понадобиться несколько десятков);
  \item Очень высокая скорость обновления системы: в течение нескольких недель после выхода новой версии её устанавливают 70–80\% пользователей, так что новыми API можно смело пользоваться, не оглядываясь на «систему пятилетней давности»;
  \item Хорошая привычка платить: пользователи App Store охотнее тратят деньги на приложения и встроенные покупки;
  \item Проверка строгая, зато магазин чистый: подделок и вредоносных программ мало, и хорошим продуктам легче быть замеченными.
\end{itemize}

\begin{tcbwarn}{Внимание: сначала о затратах}
У разработки под iOS есть жёсткий барьер: нужен Mac (Xcode существует только для macOS), а для публикации — аккаунт разработчика за 99 долларов в год (бесплатный аккаунт физлица позволяет отлаживать на реальном устройстве локально, но не публиковать приложение). Трезво оцените эти вложения перед стартом — впрочем, хорошая новость: на этапе обучения бесплатного аккаунта вполне достаточно, об этом расскажем в главе 60.

\end{tcbwarn}

\booksection{На каком языке писать: Objective-C или Swift?}

В истории разработки под iOS тоже было два «главных героя». Objective-C — ветеран-основатель: использовался с 2008 по 2014 год; Swift — новый язык Apple, представленный в 2014-м и с тех пор доросший до Swift 6 (релиз 2024 года со строгой проверкой конкурентности). К 2026 году почти все новые проекты пишутся на Swift.

Честное слово новичку, начинающему с нуля: учите сразу Swift, а Objective-C даже не открывайте — разве что на работе вам достанется поддерживать legacy-код десятилетней давности. Плюсы Swift: чистый синтаксис (никакого багажа эпохи C — точек с запятой и заголовочных файлов), механизм опциональных значений (Optional), искореняющий «падения из-за нулевых указателей» на уровне языка, бесшовная работа со SwiftUI, а ещё Swift Playgrounds на iPad — можно учиться, играя.

Стиль Objective-C (древние руины)

\begin{lstlisting}
NSString *name = @"Иван"; [name stringByAppendingString:@"Привет"]; // квадратные скобки + указатели + точка с запятой — классическая троица
\end{lstlisting}

Стиль Swift (современный и лаконичный)

\begin{lstlisting}
var name = "Иван" name += "Привет" // просто и понятно, как естественный язык
\end{lstlisting}

\booksection{Как рождается приложение?}

Как и на любой платформе, у приложения iOS есть свой конвейер. Но в версии Apple к нему добавлены две «процедуры досмотра»:

Идея

Проектирование (интерфейс + функции)

Пишем код (Swift + SwiftUI)

Сборка и подпись

Тесты (эмулятор + реальное устройство)

Проверка Apple

Публикация в App Store

\begin{center}\small\color{gray}Рисунок 1. Полный путь приложения iOS от идеи до полки магазина (обратите внимание на дополнительный этап «проверка»)\end{center}

Ключевые понятия этого процесса:

\begin{itemize}
  \item Подпись (Signing): каждое приложение iOS должно быть «подписано», чтобы его можно было установить, — это доказательство того, что «пакет исходит от вас и не был подделан». На этапе обучения Xcode автоматически подписывает приложения вашим Apple ID; для публикации используется сертификат разработчика. Подробнее — в главе 67.
  \item Симулятор (Simulator): виртуальный iPhone на Mac, главная арена разработки и отладки. Он работает на процессоре Mac, поэтому особенно плавный — по ощущениям заметно лучше эмулятора Android.
  \item TestFlight: официальное средство Apple для «раздачи бета-версий» — до публикации тестовую версию получают до 10 000 тестеров, которые устанавливают её через приложение TestFlight. Это изящное решение, свойственное именно экосистеме iOS.
\end{itemize}

\booksection{Технологический стек iOS целиком: ваш «арсенал»}

\begin{center}
\begin{tabularx}{\linewidth}{|X|X|X|}
\hline
\textbf{Уровень} & \textbf{Основной выбор (2026)} & \textbf{Для чего нужен} \tabularnewline
\hline
Язык & Swift 6 & Язык для написания логики \tabularnewline
Интерфейс & SwiftUI + дизайн-правила Apple (HIG) & Рисовать интерфейс кодом (основная линия книги) \tabularnewline
Архитектура & MVVM + @Observable & Разделять интерфейс и логику \tabularnewline
Хранение данных & SwiftData (база данных) + UserDefaults (настройки) & Сохранять данные в телефоне \tabularnewline
Сеть & URLSession + async/await + Codable & Общаться с сервером \tabularnewline
Инструмент сборки & Xcode Build System / xcodebuild & Превращать код в установочный пакет \tabularnewline
Среда разработки & Xcode (только для macOS) & Официальная IDE Apple \tabularnewline
\hline
\end{tabularx}
\end{center}

\begin{center}\small\color{gray}Обзор основного технологического стека iOS в 2026 году\end{center}

Заметили? Этот список заметно короче андроидовского — потому что Apple многое «берёт на себя». Плата за закрытую экосистему — меньше свободы, а награда — спокойствие: не нужно мучиться выбором сетевой библиотеки или базы данных; то, что даёт Apple, чаще всего и есть оптимальное решение.

\begin{tcbpractice}{Попробуйте сами: посмотрите «яблочный паспорт» своего телефона}
Откройте на iPhone [Настройки] [Основные] [Об этом устройстве] и запишите три вещи: версию iOS, название модели и доступный объём. Затем полистайте вкладку «Сегодня» в App Store и присмотритесь к скриншотам и описаниям качественных приложений — с сегодняшнего дня вы смотрите на магазин уже не глазами пользователя, а глазами разработчика.

\end{tcbpractice}

\begin{tcbsummary}{Краткие итоги главы}
\begin{itemize}
  \item В 2008 году запустился App Store — родилась «экономика приложений»; текущая версия в 2026 году — iOS 27
  \item Баланс «сада с высокими стенами»: правил много, но устройств мало, версии свежие, пользователи готовы платить, магазин чистый
  \item Жёсткие барьеры: нужен Mac; для публикации — аккаунт разработчика за 99 долларов в год; на этапе обучения хватит бесплатного аккаунта
  \item Язык — сразу Swift; Objective-C остаётся достоянием legacy-кода
  \item У конвейера на один шаг больше, чем у Android: проверка Apple; TestFlight — волшебный инструмент раздачи бета-версий
  \item Стандартный стек 2026 года: Swift 6 + SwiftUI + MVVM + SwiftData + URLSession + Xcode
\end{itemize}

\end{tcbsummary}
